
\documentclass[10pt,a4paper]{article}
\usepackage[a4paper,top=14mm,bottom=14mm,left=15mm,right=15mm,headheight=13pt]{geometry}
\usepackage{fontspec}
\usepackage{xcolor}
\usepackage{booktabs}
\usepackage{tabularx}
\usepackage{array}
\usepackage{enumitem}
\usepackage{tcolorbox}
\usepackage{fancyhdr}
\usepackage{lastpage}
\usepackage[hidelinks]{hyperref}
\usepackage{microtype}
\usepackage{amsmath}
\usepackage{titlesec}
\usepackage{ragged2e}

\setmainfont{Noto Serif}
\setsansfont{Inter}
\setmonofont{TeXGyreCursor}

\definecolor{MainBlue}{HTML}{0B2D5B}
\definecolor{SoftBlue}{HTML}{EFF5FF}
\definecolor{SoftGray}{HTML}{F6F7F9}
\definecolor{MidGray}{HTML}{D9DEE7}
\definecolor{DarkGray}{HTML}{40464F}

\pagestyle{fancy}
\fancyhf{}
\lhead{\sffamily\small Supply-Chain Knowledge Graph}
\rhead{\sffamily\small V1 Design Brief}
\cfoot{\sffamily\footnotesize \thepage/\pageref{LastPage}}
\renewcommand{\headrulewidth}{0.35pt}

\titleformat{\section}
  {\sffamily\bfseries\large\color{MainBlue}}{}{0pt}{}
\titlespacing*{\section}{0pt}{7pt}{3pt}

\setlength{\parindent}{0pt}
\setlength{\parskip}{3pt}
\setlist[itemize]{leftmargin=1.25em,topsep=2pt,itemsep=1pt,parsep=0pt}
\setlist[enumerate]{leftmargin=1.35em,topsep=2pt,itemsep=1pt,parsep=0pt}
\renewcommand{\arraystretch}{1.12}

\newtcolorbox{briefbox}[1][]{
  colback=SoftBlue,
  colframe=MainBlue,
  boxrule=0.45pt,
  arc=1.5mm,
  left=2mm,right=2mm,top=1.4mm,bottom=1.4mm,
  #1
}
\newtcolorbox{graybox}[1][]{
  colback=SoftGray,
  colframe=MidGray,
  boxrule=0.4pt,
  arc=1.5mm,
  left=2mm,right=2mm,top=1.3mm,bottom=1.3mm,
  #1
}

\begin{document}

{\sffamily\bfseries\fontsize{17}{19}\selectfont\color{MainBlue}
Supply-Chain Knowledge Graph: Frozen V1 Architecture}

\vspace{1mm}
\begin{briefbox}
\small
\textbf{Goal.} Build a reproducible supply-chain KG from unstructured text. V1 focuses on data acquisition, LLM-based entity/relation extraction, entity disambiguation, controlled upstream expansion, and Neo4j storage. Risk analysis and existing-edge monitoring are outside V1.
\end{briefbox}

\section{1. Decisions on the four open questions}

\small
\begin{tabularx}{\textwidth}{>{\bfseries\sffamily}p{0.22\textwidth} X}
\toprule
Question & Frozen V1 decision \\
\midrule
Seed definition &
Seeds are configurable and can be increased. The default replication config uses AlMahri's 13 selected companies: Toyota, Tesla, BMW, BYD; Samsung SDI, CATL, Panasonic; Zijin Mining, Albemarle, Tianqi Lithium, Ganfeng Lithium, Norilsk Nickel, Jinchuan Group. \\

Graph expansion &
Node-centric and tier-aware, not pairwise. Each frontier company is searched with a small set of tier-specific query templates. Newly discovered upstream companies become the next frontier. \\

Boundary / stopping &
\texttt{max\_tier} is configurable. Following Prof. Kuo's suggestion for testing, the default is \textbf{3 upstream supply-chain tiers} from the final EV product: seeds are Tier 0, discovered suppliers are Tiers 1--3, and Tier 3 is terminal. Only supply-chain relations (mainly \texttt{SUPPLIES\_TO}) advance the tier. \\

Ontology &
Static in V1. We use the AlMahri evaluation schema: entity types \texttt{Company, Location, Material, Person, Product, Mine}; relations \texttt{suppliesTo, locatedIn, owns, produces}. A later mapping to SCRO can be added without changing the extraction pipeline. \\
\bottomrule
\end{tabularx}

\section{2. AlMahri baseline, data modes, and datasets}

AlMahri et al. created a controlled evaluation corpus of \textbf{1,277 Wikipedia sentences} and manually annotated it for six entity types and four relation types. Their method uses zero-shot LLM prompting for NER/RE and a second prompt for entity disambiguation. The exact annotated corpus was described in the paper, but we did not identify a public downloadable package for the full labeled dataset. Because that exact labeled package is not available to us, V1 is dataset-agnostic and exposes two operating modes: \texttt{INGEST} for a fixed corpus/dataset already in hand, and \texttt{EXPAND} when the system must discover documents from seeds.

\begin{tabularx}{\textwidth}{>{\bfseries\sffamily}p{0.16\textwidth} X}
\toprule
Mode & Behavior \\
\midrule
\texttt{INGEST} &
The user provides documents or a dataset. The system normalizes, extracts facts, resolves duplicates, and merges them into Neo4j. No search is performed. This is the closest mode to the fixed-data evaluation setup in AlMahri. \\

\texttt{EXPAND} &
The user provides seeds. For each frontier company, the system runs tier-specific queries over enabled retrieval sources (e.g., Wikipedia search and Google News / GDELT), extracts upstream suppliers from the returned documents, and continues until the configured tier boundary. Wikipedia is a source in the search layer; the system does not assume that the entity's own Wikipedia page is the document to retrieve. \\
\bottomrule
\end{tabularx}

\textbf{Auxiliary reference dataset.} We use \textbf{Sectivia} as an external graph sanity-check dataset: 350 directed supplier-to-customer relations across 214 listed companies (215 company rows in the node file), released as CSV/JSON under CC BY 4.0. It is useful for checking discovered links, but it is \emph{not} a substitute for AlMahri's annotated text corpus.

\section{3. V1 architecture}

\begin{graybox}
\centering\sffamily\small
\textbf{INGEST:} Documents
$\rightarrow$ Normalize/Store
$\rightarrow$ LLM Prompt 1
$\rightarrow$ Prompt 2
$\rightarrow$ Facts
$\rightarrow$ Neo4j MERGE
\\[1mm]
\textbf{EXPAND:} Seeds
$\rightarrow$ Persistent Frontier
$\rightarrow$ Tier-Specific Retrieval Queries
$\rightarrow$ Wikipedia/Search/News Sources
$\rightarrow$ Dedup/Filter
$\rightarrow$ LLM Prompt 1
$\rightarrow$ Prompt 2
$\rightarrow$ Neo4j
$\rightarrow$ New Upstream Companies
$\rightarrow$ Next Tier
\end{graybox}

The implementation should reuse Neo4j GraphRAG KG Builder components for splitting, schema-guided entity/relation extraction, and graph writing where practical. Custom code is limited to source retrieval, expansion policy, checkpointing, provider adapters, and AlMahri-style disambiguation.

\newpage

\section{4. Graph expansion algorithm}

\begin{enumerate}
\item Initialize a persistent queue with user seeds at Tier 0.
\item Pop one \texttt{Company} and load the query templates and retrieval-source settings for its tier.
\item Run each query independently over the enabled sources (e.g., Wikipedia search, Google News, GDELT). Do not generate pairwise company queries.
\item Merge returned documents/results, normalize URLs, remove duplicates, and keep at most the configured number of results \emph{per query}. Fetch/filter only useful documents.
\item Extract facts. If a new company has a valid upstream supply relation to the current node, assign \texttt{tier=current+1}, merge it into the KG, and enqueue it only if \texttt{tier < max\_tier}. Tier-3 nodes are therefore stored but not expanded under the default setting.
\item Continue until the frontier is empty or the tier/node limits are reached.
\end{enumerate}

\begin{briefbox}
\small
\textbf{Query design.} We never create all company pairs. Queries are \textbf{node-centric}: each tier has a configurable set of keywords/query templates, and the current frontier company's name is inserted into them. Search cost is approximately $O(NK)$, where $N$ is the number of expanded companies and $K$ is the small number of templates for that tier, rather than $O(N^2)$.

\textbf{Example.} If BMW is a Tier-0 frontier node, Tier-0 templates such as \texttt{"\{entity\}" suppliers}, \texttt{"\{entity\}" sourcing}, and \texttt{"\{entity\}" battery suppliers} become \texttt{"BMW" suppliers}, \texttt{"BMW" sourcing}, and \texttt{"BMW" battery suppliers}. If CATL is then discovered as an upstream supplier, it becomes a Tier-1 frontier node and is searched with the Tier-1 set, e.g., \texttt{"CATL" suppliers}, \texttt{"CATL" component suppliers}, and \texttt{"CATL" raw material suppliers}. Tier-2 templates can focus further upstream, e.g., \texttt{"\{entity\}" material suppliers} and \texttt{"\{entity\}" mines}. Each template can have its own result limit through the same config hierarchy described below.
\end{briefbox}

\textbf{Config inheritance.} Search limits use ``closest value wins'':
\[
\text{global defaults} \;<\; \text{tier defaults} \;<\; \text{query override}.
\]
For example, the global default may be 20 results/query, Tier 0 may override it to 30, and one specific query may override it again to 15. This rule is applied per field, so only explicitly changed values are replaced.

\section{5. LLM prompts and extraction}

\begin{tabularx}{\textwidth}{>{\bfseries\sffamily}p{0.18\textwidth} X}
\toprule
Prompt & V1 behavior \\
\midrule
Prompt 1: NER + RE &
\textbf{Purpose: extract graph facts from text.} The prompt receives a document/chunk plus the fixed ontology and asks the model to return only allowed entity types and relations. Its output is structured entities and triples, together with the supporting text/source evidence. For example, from ``CATL supplies battery cells to BMW'', it should identify CATL and BMW as \texttt{Company} entities and produce \texttt{CATL suppliesTo BMW}. NER and relation extraction are conceptually separate tasks, but V1 can obtain both in one structured LLM call. \\

Prompt 2: Entity disambiguation &
\textbf{Purpose: merge mentions that refer to the same real-world entity.} After Prompt 1, mentions are grouped by entity type and passed to this second prompt. The model assigns the same canonical/group ID to equivalent mentions, e.g., ``CATL'' and ``Contemporary Amperex Technology Co. Limited'', while keeping genuinely different entities separate. The resolved IDs are then used before Neo4j \texttt{MERGE}, preventing duplicate nodes. \\
\bottomrule
\end{tabularx}

\begin{graybox}
\small
\textbf{Prompt flow:} document text + static ontology $\rightarrow$ \textbf{Prompt 1} $\rightarrow$ extracted entities/relations/evidence $\rightarrow$ \textbf{Prompt 2} $\rightarrow$ canonical entity groups $\rightarrow$ deduplicated facts $\rightarrow$ Neo4j.
\end{graybox}

The LLM layer is provider-independent. The planned local model is \textbf{Qwen3.8-27B} in instruct/non-thinking mode, served with vLLM; a GPT API can be kept as an optional comparison/provider. The same extraction contract is used in either case, so the KG logic does not depend on the model backend.

\section{6. Reliability and scope}

\begin{itemize}
\item \textbf{Resume is default.} Both modes save state in a small persistent store (e.g., SQLite): document status for \texttt{INGEST}, and frontier/search/document status for \texttt{EXPAND}. A user can explicitly request a fresh run.
\item \textbf{No duplicate work.} Raw documents are stored with normalized URL/content hash; normalized data and extracted facts are saved separately (JSONL/Parquet). Completed LLM results are reused after restart.
\item \textbf{Idempotent graph writes.} Canonical entity IDs plus Neo4j \texttt{MERGE} prevent duplicate nodes/edges when a run is resumed or repeated.
\item \textbf{Not in V1.} Existing-edge monitoring, relationship-change detection, risk analytics, and custom graph dashboards. Neo4j Browser is sufficient for inspection/visualization.
\end{itemize}

\vspace{1mm}
{\footnotesize
\textbf{References.}
[1] S. AlMahri, L. Xu, A. Brintrup, ``Enhancing Supply Chain Visibility with Knowledge Graphs and Large Language Models,'' IJPR / arXiv:2408.07705.
[2] Neo4j GraphRAG Python, Knowledge Graph Builder documentation.
[3] IOFoundry, Supply Chain Reference Ontology (SCRO).
[4] Sectivia, Open Supply-Chain Dataset (CC BY 4.0).
}

\end{document}
